% Experiments
\section{Experiments}
\label{sec:experiments}

Our experimental evaluation focuses on characterizing the performance of \texttt{gce\_tangent} against \texttt{gce\_ambient} across three distinct concept domains. We prioritize nonlinear concepts where manifold geometry is expected to play a critical role, while maintaining a diagnostic baseline for linear comparison.

\subsection{Experimental Domains and Ordering}
We evaluate our methods across three primary blocks, ordered by their relevance to nonlinear manifold erasure:
\begin{enumerate}
    \item \textbf{Sycophancy (Primary Nonlinear Case):} Our central focus is on removing agreement bias in model responses. This represents a critical challenge for model alignment where concept encoding is expected to be locally curved.
    \item \textbf{Safety (Secondary Case Study):} We evaluate the surgicality of erasure on content safety using the PKU-SafeRLHF dataset, testing the preservation of helpfulness alongside safety leakage.
    \item \textbf{Bias in Bios (Diagnostic Baseline):} We use gender and profession classification as a replication anchor. This setting provides a linear comparison point to contrast against the nonlinear gains of tangent-constrained erasure.
\end{enumerate}

\subsection{Model and Implementation}
All experiments are conducted on Qwen2.5-0.5B (24 decoder layers, hidden dimension 896). The current paper-locked comparison set implements five primary erasure methods: linear projection, INLP, LEACE, \texttt{gce\_ambient}, and \texttt{gce\_tangent}. LEACE is included as a matched linear baseline in the layer-8 inventory reported below. For \texttt{gce\_tangent}, we estimate local tangent spaces with PCA-based neighborhoods so that the erasure direction is constrained to a local on-manifold approximation rather than to the full ambient space.

\begin{table}[t]
\centering
\caption{Comparison of concept erasure methods by family and role in this work.}
\label{tab:method_comparison_scope}
\scriptsize
\begin{tabular}{llp{0.48\linewidth}}
\toprule
Method & Family & Role in Paper \\
\midrule
Linear Projection & Linear & Global baseline \\
INLP & Linear & Iterative linear baseline \\
LEACE & Linear & Strong matched linear baseline \\
IGBP & Nonlinear & Closest prior nonlinear baseline \\
\AmbGCE{} & Nonlinear & Ambient ablation of \TanGCE{} \\
\TanGCE{} & Nonlinear & Main contribution \\
\bottomrule
\end{tabular}
\end{table}


\subsection{Evaluation Protocol}
All leakage probes are trained from scratch on erased representations to ensure no information leaks from the erasure scorers themselves. We employ three complementary evaluation axes:

\paragraph{Probe-based leakage.} A retrained nonlinear MLP probe measures whether concept information remains accessible to a standard function class. We report both linear and nonlinear test accuracy; the gap between them indicates the degree to which concept encoding is nonlinear.

\paragraph{Mutual information via $k$-NN estimation.}
Probe accuracy measures whether a \emph{specific function class} can recover the concept; a nonlinear MLP probe might miss information that a more complex architecture could extract. Mutual information $I(\mathbf{X}; y)$ captures \emph{all} shared information between representations and labels, providing a probe-free, architecture-agnostic bound on recoverability.

\emph{Mental model.} Mutual information decomposes as $I(\mathbf{X}; y) = H(y) - H(y \mid \mathbf{X})$: the entropy of the concept label minus the residual uncertainty after observing the representation. When erasure is effective, $H(y \mid \mathbf{X})$ rises toward $H(y)$ (the label becomes unpredictable from the representation), and MI drops toward zero. For continuous $\mathbf{X}$ and discrete $y$, we estimate MI using the Ross/Gao $k$-nearest-neighbor estimator:
\begin{equation}
\label{eq:knn_mi}
I(\mathbf{X}; y) = \psi(k) - \bigl\langle \psi(m_i) \bigr\rangle + \psi(N) - \bigl\langle \psi(N_{c(i)}) \bigr\rangle,
\end{equation}
where $\psi$ is the digamma function and the procedure is:
\begin{enumerate}
    \item For each point $\mathbf{x}_i$, find its $k$-th nearest neighbor \emph{within its class} $c(i)$; call that distance $\varepsilon_i$.
    \item Count $m_i$: the number of \emph{all} points (any class) within distance $\varepsilon_i$ of $\mathbf{x}_i$.
    \item $N_{c(i)}$ is the total number of points sharing the same label as point $i$; $N$ is the total sample size.
\end{enumerate}

\emph{Intuition.} When concept information is encoded in $\mathbf{X}$, same-class points cluster together: $\varepsilon_i$ is small, $m_i \approx k$ (mostly same-class neighbors within that radius), and MI is high. After effective erasure, class labels become geometrically mixed: to reach $k$ same-class neighbors one must search a large $\varepsilon_i$, sweeping in many cross-class points so that $m_i \gg k$, and MI drops.

\emph{Adaptive resolution and local geometry.} A key property of the $k$-NN estimator is that it sets its measurement scale \emph{adaptively}: for each point, the radius $\varepsilon_i$ is determined by the local data density, not by a global bandwidth or histogram bin width. In dense regions of the manifold, $\varepsilon_i$ is small; in sparse regions, it is large. The digamma terms $\psi(m_i)$ implicitly encode the local density ratio between the class-conditional and marginal distributions without ever estimating densities explicitly in $\mathbb{R}^d$---the estimator works entirely with neighbor counts, which are topological quantities that naturally handle varying density and curvature.

This makes the estimator well-suited for data on a low-dimensional manifold embedded in $\mathbb{R}^d$. A parametric density estimator (e.g., Gaussian kernel) operating in the ambient space would be misled by the high-dimensional volume: two points close in $\mathbb{R}^d$ may be far apart along the manifold, and vice versa. The $k$-NN estimator sidesteps this because it only examines local neighborhoods---points that are actually nearby along the data distribution. While it still uses ambient Euclidean distance (not geodesic distance), the locality of $k$-NN queries means that for small $k$ relative to the dataset, the Euclidean and geodesic neighborhoods approximately coincide.

\emph{Connection to TanGCE.} Both the MI estimator and TanGCE operate through local neighborhood structure rather than global ambient-space assumptions. TanGCE explicitly projects onto the tangent space estimated from $k$-nearest neighbors; the MI estimator implicitly adapts to the manifold by measuring information through local ball neighborhoods. Neither assumes the data fills $\mathbb{R}^d$ uniformly. This shared reliance on local structure creates a coherent geometric framework: the same manifold properties that make tangent-projected erasure more precise also make $k$-NN MI estimation more faithful than ambient-space alternatives.

We compute MI for both concept labels (should decrease after erasure) and control labels (should remain unchanged), yielding a complete picture of erasure selectivity.

\paragraph{Unsupervised SAE diagnostics.} We train TopK sparse autoencoders ($k{=}64$) on a general-purpose WikiText corpus---independent of any concept labels---then measure concept feature removal rate, collateral damage rate, and feature redistribution score (Section~\ref{sec:results}). These metrics require no control-task annotations and detect failure modes invisible to probe-based evaluation, such as concept information being redistributed across features rather than removed.

\subsection{Geometric Foundation (Appendix Support)}
While the main text focuses on the performance of \texttt{gce\_tangent}, we also provide geometric diagnostics to ground the mechanism story. These include broad geometry sweeps across multiple datasets and layers, subspace projection sensitivity (Q1), perturbation-based subspace sensitivity tests (Q2), and representation connectivity analyses (Q3). These appendix-scale materials support the low-dimensional and nonlinear-geometry motivation for tangent projection, but they should not be read as a standalone causal proof that off-manifold displacement fully explains the observed utility tradeoffs.
